\documentclass{ximera}
\typeout{Start loading xmPreamble.tex}%

% Add here extra macro's that are loaded automatically by all documents of claas 'ximera' or 'xourse' in this repo

%%
%%  Example:
%%
% \newcommand{\R}{\mathbb{R}}
\usepackage{tikz}
\usetikzlibrary{positioning, arrows.meta}
\usepackage{multicol}
\usepackage{enumitem}
\usepackage{caption}
\usepackage{amsmath}
\usepackage{amsthm}

\title{Classification of States}
\begin{document}
\begin{abstract}
An introduction to the classification of states in Markov chains and processes, covering communicating classes, irreducibility, recurrence, transience, stationary distributions, reversibility, and an application to the Asymmetric Simple Exclusion Process.
\end{abstract}
\maketitle

\section*{Classification of States}

\begin{definition}
For $x, y \in \mathcal{X}$, we write $x \to y$ ($y$ can be reached from $x$) if $P^n_{xy} > 0$ for some $n \geq 0$. In the continuous-time setting, this would be $P_{xy}(t) > 0$ for some $t \geq 0$. Furthermore, write $x \leftrightarrow y$ ($x$ and $y$ communicate) if $x \to y$ and $y \to x$.
\end{definition}

\begin{proposition}
The relation $\leftrightarrow$ defines an equivalence relation.
\end{proposition}

\begin{proof}
Observe that $x \leftrightarrow x$ (take $n = 0$), so the reflexive property holds. Also, $x \leftrightarrow y$ if and only if $y \leftrightarrow x$, so the symmetric property holds. To prove transitivity, note that for a Markov chain,
\[
P^{n+m}_{xz} \geq P^n_{xy} P^m_{yz}
\]
while for a Markov process,
\[
P_{xz}(s+t) = P_{xy}(s)P_{yz}(t).
\]
So $x \leftrightarrow y$ and $y \leftrightarrow z$ imply $x \leftrightarrow z$.
\end{proof}

\begin{definition}
A \textbf{communicating class} $C \subseteq \mathcal{X}$ is an equivalence class of $\leftrightarrow$. So in particular, $x \leftrightarrow y$ for all $x, y \in C$. Furthermore, $\mathcal{X}$ is \textbf{irreducible} if $\mathcal{X}$ itself is a communicating class, i.e.\ $x \leftrightarrow y$ for all $x, y \in \mathcal{X}$.
\end{definition}

In general, $\mathcal{X}$ is a disjoint union of communicating classes.

\begin{example}
Consider the matrix
\[
P = \begin{pmatrix} 1/2 & 1/2 & 0 \\ 1/4 & 1/2 & 1/4 \\ 0 & 1/3 & 2/3 \end{pmatrix}
\]
with rows and columns labeled $1, 2, 3$. Then $1 \to 2$, $2 \to 1$, $2 \to 3$, $3 \to 2$, which gives $1 \leftrightarrow 2$ and $2 \leftrightarrow 3$. So $\{1, 2, 3\}$ is a single communicating class and the chain is irreducible.
\end{example}

\begin{example}
Consider the matrix
\[
P = \begin{pmatrix} 1/2 & 1/2 & 0 & 0 \\ 1/2 & 1/2 & 0 & 0 \\ 0 & 0 & 1/4 & 3/4 \\ 0 & 0 & 0 & 1 \end{pmatrix}
\]
Then $1 \leftrightarrow 2$ so $\{1,2\}$ is a communicating class. We also have $3 \to 4$ but $4 \not\to 3$, so $\{3\}$ and $\{4\}$ are each their own communicating classes. In general, if $P$ is a block matrix of the form
\[
P = \begin{pmatrix} A & 0 \\ 0 & B \end{pmatrix}
\]
then the Markov chain is not irreducible.
\end{example}

\begin{definition}
(Informal) A state is \textbf{transient} if there is a positive probability of never returning to it. A state is \textbf{recurrent} if it is not transient.
\end{definition}

\begin{definition}
Let $\tau_x = \inf\{n \geq 1 : X_n = x\}$, called the \textbf{first hitting time} of $x$. Let $f_x = P(\tau_x < \infty \mid X_0 = x)$. A state is \textbf{transient} if $f_x < 1$ and \textbf{recurrent} if $f_x = 1$.
\end{definition}

\begin{theorem}
\begin{enumerate}[label=(\alph*)]
\item Suppose $X_0 = x$. Then the number of visits to $x$ is distributed as a geometric random variable with parameter $1 - f_x$.
\item $x$ is recurrent if and only if $\displaystyle\sum_{n=1}^{\infty} P^n_{xx} = \infty$.
\end{enumerate}
\end{theorem}

\begin{proof}
(a) After each visit to $x$,
\[
P(\text{never returning to } x) = P(\tau_x = \infty \mid X_0 = x) = 1 - f_x.
\]
Meanwhile,
\[
P(\text{exactly } n \text{ visits}) = f_x^n(1 - f_x)
\]
because each visit occurs independently with probability $f_x$. So the number of visits is distributed as $\text{Geo}(1 - f_x)$.

(b) By (a) and the expectation of a geometric random variable,
\[
E_x[\text{number of visits to } x] = \frac{1}{1 - f_x}.
\]
Calculating the same quantity another way,
\[
E_x[\#\text{ visits to } x] = E\left[\sum_{n=1}^{\infty} \mathbf{1}_{\{X(n)=x\}}\right] = \sum_{n=1}^{\infty} P_x(X(n) = x) = \sum_{n=1}^{\infty} P^n_{xx}.
\]
Therefore $x$ is recurrent $\iff$ $f_x = 1$ $\iff$ $\displaystyle\sum_{n=1}^{\infty} P^n_{xx} = \infty$.
\end{proof}

\begin{theorem}
A communicating class consists entirely of recurrent states or entirely of transient states.
\end{theorem}

\begin{proof}
Suppose $x$ is recurrent and $x \leftrightarrow y$. Then $P^k_{xy} > 0$ and $P^l_{yx} > 0$ for some $k, l \geq 0$. Using
\[
P^{l+n+k}_{yy} \geq P^l_{yx} P^n_{xx} P^k_{xy},
\]
we conclude that $\displaystyle\sum_{n=0}^{\infty} P^{l+n+k}_{yy} = \infty$ because $x$ is recurrent. Therefore $\displaystyle\sum_{m=0}^{\infty} P^m_{yy} = \infty$ and $y$ is also recurrent.
\end{proof}

\begin{remark}
If the state space $\mathcal{X}$ is finite, then $\mathcal{X}$ has at least one recurrent state, since otherwise $X_n$ would have nowhere to go.
\end{remark}

\begin{example}
Consider the infinite state space $\mathcal{X} = \mathbb{Z}$ with $X(0) = 0$ and an asymmetric simple random walk where right jumps have probability $p$ and left jumps have probability $1-p$. Since $\mathcal{X}$ is irreducible, we only need to check if $0$ is recurrent. For even time steps,
\[
P^{2n}_{00} = \binom{2n}{n} p^n(1-p)^n.
\]
Using Stirling's formula $n! \sim n^n e^{-n} \sqrt{2\pi n}$, we get $\binom{2n}{n} \sim \frac{2^{2n}}{\sqrt{\pi n}}$, so
\[
P^{2n}_{00} \sim \frac{[4p(1-p)]^n}{\sqrt{\pi n}}.
\]
If $p = 1/2$, then $P^{2n}_{00} \sim \frac{1}{\sqrt{\pi n}}$ and $\sum P^{2n}_{00} = \infty$, so $0$ is recurrent. If $p \neq 1/2$, then $4p(1-p) < 1$ and $\sum P^{2n}_{00} < \infty$, so $0$ is transient. Thus a symmetric simple random walk is recurrent and an asymmetric one is transient.
\end{example}

\begin{definition}
A recurrent state $x$ is \textbf{positive recurrent} if $E_x[\tau_x] < \infty$, and \textbf{null recurrent} if $E_x[\tau_x] = \infty$.
\end{definition}

\begin{theorem}
Suppose $x \leftrightarrow y$, both recurrent. If $y$ is positive recurrent, then so is $x$.
\end{theorem}

\begin{proof}
Let $\tau_{xy} = \min_{n \geq 1}\{X(n) = y \mid X(0) = x\}$ be the first time to go from $x$ to $y$. Since $\tau_{xx} \leq \tau_{xy} + \tau_{yx}$, it suffices to show $E[\tau_{xy}], E[\tau_{yx}] < \infty$.

For $E[\tau_{xy}] < \infty$: since $x \leftrightarrow y$, there is a smallest $n$ such that $P^{(n)}_{yx} > 0$. Define the event
\[
A = \{X(k) \neq y \text{ for } 1 \leq k < n \text{ and } X(n) = x\},
\]
so $P(A) > 0$. Since $y$ is positive recurrent, by the total probability rule and the Markov property,
\[
\infty > E_y[\tau_y] \geq E_y[\tau_y \mid A] P(A) = (n - 1 + E[\tau_{xy}]) P(A),
\]
and since $P(A) > 0$ we conclude $E[\tau_{xy}] < \infty$.

For $E[\tau_{yx}] < \infty$: let $Y_m$ be the time between the $m$th and $(m+1)$th visit to $y$. By the Markov property, $Y_m$ are i.i.d.\ with $E[Y_m] = E[\tau_y] < \infty$. Let $p = P_y(X(n) \text{ visits } x \text{ before returning to } y) \geq P(A) > 0$, and let $N \sim \text{Geo}(p)$ be the number of revisits to $y$ before first visiting $x$. Then $\tau_{yx} \leq Y_1 + \cdots + Y_{N+1}$, so $E[\tau_{yx}] \leq E[Y_0] E[N] < \infty$.
\end{proof}

For continuous-time processes:

\begin{definition}
A state $x$ is:
\begin{enumerate}[label=(\alph*)]
\item \textbf{recurrent} if $\displaystyle\int_0^{\infty} P_{xx}(t)\, dt = \infty$,
\item \textbf{transient} if $\displaystyle\int_0^{\infty} P_{xx}(t)\, dt < \infty$,
\item \textbf{positive recurrent} if $\lim_{t \to 0} P_{xx}(t) > 0$, and \textbf{null recurrent} if $\lim_{t \to 0} P_{xx}(t) = 0$.
\end{enumerate}
These are again class properties.
\end{definition}

\section*{Stationary States and Reversibility}

\begin{definition}
For a state $y \in \mathcal{X}$, let $\pi_y$ be the long-run proportion of time spent at $y$:
\[
\pi_y = \lim_{n \to \infty} \frac{1}{n} \sum_{m=1}^{n} \mathbf{1}_{\{X(m)=y\}},
\]
where this is an almost sure limit. If $\pi_y$ exists and does not depend on the initial state, we call $\pi$ the \textbf{limiting} or \textbf{stationary distribution}.
\end{definition}

\begin{proposition}
If a Markov chain on an irreducible state space is positive recurrent, then there is a unique stationary distribution given by
\[
\pi_x = \frac{1}{E[\tau_x]} > 0 \quad \forall x.
\]
\end{proposition}

\begin{proof}
Let $Y_1, \ldots, Y_k$ be the times between revisits to $x$. The $k$th visit occurs at time $Y_1 + \cdots + Y_k$, so
\[
\lim_{n \to \infty} \frac{1}{n} \sum_{m=1}^{n} \mathbf{1}_{\{X_n = x\}} = \lim_{k \to \infty} \frac{k}{Y_1 + \cdots + Y_k}.
\]
Since $Y_1, \ldots, Y_k$ are i.i.d.\ with mean $E[\tau_x] < \infty$, by the strong law of large numbers $\pi_x = \frac{1}{E[\tau_x]}$.
\end{proof}

\begin{proposition}
If $X(n)$ is positive recurrent and irreducible with stationary distribution $\pi$, then $\pi P = \pi$.
\end{proposition}

\begin{proof}
Starting from $\lim_{n\to\infty} \frac{1}{n} \sum_{m=1}^n P^m = \pi$ (as a matrix with each row equal to $\pi$), multiplying both sides by $P$ and using the fact that $P^k$ is stochastic (entries bounded between 0 and 1) gives $\pi P = \pi$.
\end{proof}

\begin{definition}
Given an irreducible Markov chain $X(n)$ with transition matrix $P$ and stationary measure $\pi$, define the \textbf{reversed chain} by
\[
P^{\text{rev}}(y, x) = \frac{\pi(y)}{\pi(x)} P(x, y).
\]
If $P^{\text{rev}} = P$, we say the Markov chain is \textbf{reversible} and call $\pi$ a \textbf{reversible measure}.
\end{definition}

Reversibility can be expressed as the intertwining relation $P^T = V^{-1} P V$, where $V$ is the diagonal matrix with entries $V_{xx} = \pi(x)$.

\section*{Application: Asymmetric Simple Exclusion Process (ASEP)}

We apply the above theory to the Asymmetric Simple Exclusion Process (ASEP) on the integer lattice with step initial conditions, where particles start to the right of the origin.

\begin{itemize}
\item For $q < 1$: particles drift left, away from their initial conditions. ASEP is \textbf{transient}.
\item For $q > 1$: particles drift right, back into where they started. ASEP is \textbf{positive recurrent}, and its reversible measures are called \textbf{blocking measures}.
\item For $q = 1$: to be determined.
\end{itemize}

On a finite lattice with a fixed number of particles, there is a unique reversible measure on each communicating class. As the lattice size tends to infinity, these reversible measures converge to the blocking measures of the infinite system.

\end{document}